\chapter{Modern molecular programming and evidence}
\label{ch:dna31}
\section{The engineering question}
A formal network can be mathematically correct while its DNA realization runs
too slowly, consumes an unreported auxiliary species, or activates a reporter
through an unintended route. These are statements about different objects,
not contradictions. This chapter connects those objects without erasing their
differences.

Chapter 30 separated genomic organization, regulation and development. Here we
return to synthetic molecular programming: what supports the assertion that a
molecular system implements a specified computation? We build an original
finite-fuel reaction model, derive an exact constant-fuel reference, implement a
positivity-preserving update, and carry uncertainty through readout and evidence
assessment. The model is small enough to audit. It is not a complete DNA compiler,
experimental data, or a wet-laboratory protocol.

The reader will distinguish formal species from physical intermediates,
calculate failures of a reduced model, diagnose a non-identifiable readout, and
reject unsupported evidence transfer. Computed tables come from
\texttt{book/chapters/ch31/reference.py}.
\Plate{01-stack}{Four interfaces in molecular programming. Success at one does not prove the others.}{fig:dna31-stack}

\section{Specify the refinement obligation}
Let $x(t)\in\mathbb R_{\ge0}^{n}$ denote formal concentrations with
$\dot x=f(x)$. A physical state $z(t)\in\mathbb R_{\ge0}^{m}$ includes auxiliary
complexes, intermediates and waste, with $\dot z=g(z;\theta)$. A
\Term{refinement map} $P$ identifies the physical representation of formal
species. A useful statement names a norm, horizon, initial domain and tolerance:
\begin{equation}
 \sup_{0\le t\le T}\|Pz(t)-x(t)\|\le\varepsilon
 \quad\text{for }z(0)\in\mathcal D,\ \theta\in\Theta.
 \label{eq:dna31-refinement}
\end{equation}
A bound for one minute need not persist for one day. A low-signal approximation
need not hold at arbitrary concentration. A correct terminal Boolean answer
does not establish trajectory refinement.

Soloveichik, Seelig and Winfree develop strand-displacement constructions for
approximating prescribed chemical kinetics \cite{Soloveichik2010}. This motivates
an explicit refinement interface. It does not identify our original, lumped
three-reaction model with their compiler, or provide parameters for our examples.

\section{A displacement mechanism is not a reaction arrow}
An exposed toehold permits initial association. Branch migration exchanges which
strand occupies a complementary domain. Release leaves the invader paired and
the incumbent free. Figure~\ref{fig:dna31-dsd} labels antiparallel polarities;
$a^*$ denotes the complementary domain, not an arbitrary second sequence.
Association, migration and dissociation are distinct microscopic events. The
figure does not imply deterministic motion, successful release after every
encounter, or experimentally validated sequences. Coarse-graining these events
into an effective rate is an assumption to test.
\Plate{02-dsd}{Original domain-level displacement schematic. Arrowheads mark $3'$ ends, including the left-pointing scaffold. Migration is reversible; positions are schematic, not time-resolved.}{fig:dna31-dsd}

\section{Build a finite-fuel realization}
Consider the well-mixed teaching network
\begin{equation}
 A+F\xrightarrow{k_{\rm on}}I,\qquad
 I\xrightarrow{k_{\rm off}}A+F,\qquad
 I\xrightarrow{k_{\rm cat}}B+W.
 \label{eq:dna31-network}
\end{equation}
$A$ is available input, $F$ auxiliary fuel, $I$ a captured intermediate,
$B$ output and $W$ spent auxiliary material. The symbol $k_{\rm cat}$ labels a
unimolecular conversion; it does not make $A$ a catalyst or imply a protein
enzyme. $k_{\rm on}$ has units concentration$^{-1}$time$^{-1}$, while the other
rates have units time$^{-1}$.

Set $v_1=k_{\rm on}AF$, $v_2=k_{\rm off}I$, $v_3=k_{\rm cat}I$.
In order $(A,F,I,B,W)$:
\begin{equation}
 \dot z=
 \begin{pmatrix}-1&1&0\\-1&1&0\\1&-1&-1\\0&0&1\\0&0&1\end{pmatrix}
 \begin{pmatrix}v_1\\v_2\\v_3\end{pmatrix}.
 \label{eq:dna31-stoich}
\end{equation}
Capture consumes both available pools, reversal restores them, and conversion
produces one output and one waste. Adding rows proves the conserved pools
\begin{equation}
 C_A=A+I+B,\qquad C_F=F+I+W.
 \label{eq:dna31-pools}
\end{equation}
$B-W$ is also constant. Initially absent $I,B,W$ imply $B=W$ and
$0\le B(t)\le\min(A_0,F_0)$. Waiting cannot overcome this capacity limit.

If formal input means all unconverted material, its projection is $X=A+I$,
not free $A$. Then $\dot X=-k_{\rm cat}I$. Comparing free concentration in one
model against total unconverted concentration in another confuses definition
error with approximation error. We retain both $A$ and $I$ below.

\section{Derive the held-fuel reference}
An external reservoir maintaining $F=\bar F$ gives constant
$u=k_{\rm on}\bar F$. This is not a sealed mixture:
\[
 \frac{d}{dt}\begin{pmatrix}A\\I\end{pmatrix}
 =
 \begin{pmatrix}-u&k_{\rm off}\\u&-(k_{\rm off}+k_{\rm cat})\end{pmatrix}
 \begin{pmatrix}A\\I\end{pmatrix}.
\]
The poles solve
$\lambda^2+s\lambda+u k_{\rm cat}=0$, where
$s=u+k_{\rm off}+k_{\rm cat}$. Put
$\Delta=\sqrt{s^2-4u k_{\rm cat}}$. For $\Delta>0$,
\[
 \lambda_f=-(s+\Delta)/2,\qquad
 \lambda_s=-2u k_{\rm cat}/(s+\Delta).
\]
The second formula avoids cancellation when the slow pole is close to zero.
Nonnegative rates give a nonnegative discriminant:
$(u-k_{\rm cat})^2+k_{\rm off}^2+2k_{\rm off}(u+k_{\rm cat})$.
For $A(0)=A_0$, $I(0)=B(0)=0$:
\begin{align}
 I(t)&=\frac{uA_0}{\Delta}(e^{\lambda_s t}-e^{\lambda_f t}),\\
 A(t)&=\frac{A_0}{\Delta}
 [(-\lambda_f-u)e^{\lambda_s t}+(\lambda_s+u)e^{\lambda_f t}],\\
 B(t)&=A_0-A(t)-I(t).
\end{align}
A repeated pole requires its limiting solution, not an arbitrary small
denominator. For $k_{\rm off}=0$, $u=k_{\rm cat}>0$:
$A=A_0e^{-ut}$, $I=A_0ut e^{-ut}$. For $u=0$, no capture occurs.
\Listing{held_fuel}

\section{Expose the reduction's initial-layer error}
Setting $\dot I\approx0$ yields
\[
 I\approx \frac{uA}{k_{\rm off}+k_{\rm cat}},\qquad
 k_{\rm eff}=\frac{u k_{\rm cat}}{k_{\rm off}+k_{\rm cat}}.
\]
Using $A\approx A_0-B$ adds a second assumption: little input resides in
$I$. A useful ratio is $u/(k_{\rm off}+k_{\rm cat})$; its adequacy depends on
the required tolerance. The reduced curve is
$B_{\rm red}=A_0(1-e^{-k_{\rm eff}t})$.

The exact initial slope is $\dot B(0)=k_{\rm cat}I(0)=0$. The reduced initial
slope is $k_{\rm eff}A_0>0$. Hence the reduction is not an exact transient
identity, even with perfect held fuel. It discards intermediate accumulation.
For $u=.1$, $k_{\rm off}=.9$, $k_{\rm cat}=.1$, $k_{\rm eff}=.01$;
the exact slow decay magnitude is approximately $.00917$.
These are assigned model units, not measured rates.
\begin{center}\input{\Generated/kinetics.tex}\end{center}
\Plate{03-crn}{Computed held-fuel comparison: exact output, reduced output and intermediate. Finite-fuel pool conservation is not imposed on an externally replenished reservoir.}{fig:dna31-crn}

\section{Implement without clipping the answer}
Use simultaneous Euler updates $z_{n+1}=z_n+hSv(z_n)$. A sufficient positivity
condition is
\begin{equation}
 h\max\{k_{\rm on}F_n,k_{\rm on}A_n,k_{\rm off}+k_{\rm cat}\}\le1.
 \label{eq:dna31-positive}
\end{equation}
For example, $A_{n+1}=A_n(1-hk_{\rm on}F_n)+hk_{\rm off}I_n$ has two
nonnegative terms. The same reasoning applies to $F$ and $I$; $B,W$ increase.
The code refuses an unsafe step rather than clamping negative concentrations.
\Listing{euler_step}

The conserved-pool vectors annihilate each stoichiometric column, so Euler
preserves those pools algebraically. Tests allow explicit roundoff tolerances.
Clipping would change a pool without a modeled reaction and conceal an
integration failure. Positivity alone is not an accuracy theorem: we test
first-order step refinement against an independently handwritten RK4 system.
The exact held-fuel formula is separately compared with a handwritten linear
integrator; neither oracle calls the production right-hand side.

The sealed-mixture table uses $(k_{\rm on},k_{\rm off},k_{\rm cat})=(.1,.9,.1)$,
$A_0=1$, other signal pools initially empty, horizon $20$, and step $.002$:
\begin{center}\input{\Generated/fuel.tex}\end{center}
Changing $F_0$ changes both capacity and initial capture rate. It does not isolate
a pure capacity effect, nor constitute a measured dose-response experiment.

\section{Composition changes the interface}
A signal captured by two approximately constant downstream populations has
hazards $\lambda_1=k_1C_1$ and $\lambda_2=k_2C_2$. Under independent exponential
waiting times,
\[
 P(\text{branch 1 wins})=
 \int_0^\infty\lambda_1e^{-(\lambda_1+\lambda_2)t}\,dt
 =\frac{\lambda_1}{\lambda_1+\lambda_2}.
\]
Tested alone, branch 1 captures with probability one if $\lambda_1>0$.
Adding equal competition makes that probability one half.
This \Term{loading counterexample} excludes dissociation, replenishment and
spatial transport. Those require a richer model, not an informal correction.
\Plate{04-composition}{Shared-signal competition. Isolated module tests do not imply composed behavior. The perturbation envelope below is a distinct, conditional bound.}{fig:dna31-composition}

\subsection{Propagate a bounded discrepancy}
Let initial separation be $E_0$, vector-field discrepancy at most $\eta$ and
Lipschitz constant $L$ on a declared domain.
$E(t)\le E_0+\int_0^t(LE(s)+\eta)\,ds$ gives
\[
 E(t)\le E_0e^{Lt}+
 \begin{cases}\eta(e^{Lt}-1)/L,&L>0,\\\eta t,&L=0.\end{cases}
\]
The implementation uses \texttt{expm1} near zero. It evaluates supplied constants;
it does not certify them. A small fitted observation residual cannot bound
unobserved vector-field discrepancy. An envelope can be finite yet too large to
support the advertised tolerance.

\section{Observe a mechanism, not merely brightness}
If a channel measures $y=B+I$, the states $(B,I)=(.4,.1)$ and $(.2,.3)$ are
indistinguishable. Lower instrument noise cannot remove this
\Term{structural non-identifiability}. An independently calibrated channel for
$I$ would resolve $B=y-I$ under that model; a conserved input total alone would
not separate $B$ from $I$.

A reporter can also brighten when degradation separates fluorophore and
quencher rather than when intended displacement occurs.
\Plate{05-environment}{Distinct intended and unintended activation routes. The cell and compartment boundaries are explanatory schematics, not microscopy or validated delivery models.}{fig:dna31-environment}
Fern and Schulman examine nuclease-associated degradation in serum-supplemented
medium \cite{FernSchulman2017}. Zhong and Sczepanski compare D- and L-DNA displacement
behavior in living mammalian cells \cite{ZhongSczepanski2021}. These primary studies
motivate context-specific records; their rates are not parameters of our model.

\section{Carry calibration uncertainty before thresholding}
For nonsaturating $y=o+gB$, assume independent interval bounds
$y\in[y_-,y_+]$, $o\in[o_-,o_+]$, $0<g_-\le g\le g_+$, $B\ge0$.
A feasible enclosure is
\begin{equation}
 B\in\left[\max\left(0,\frac{y_--o_+}{g_+}\right),
                 \frac{y_+-o_-}{g_-}\right].
 \label{eq:dna31-calibration}
\end{equation}
If $y_+<o_-$, nonnegative output cannot explain the observation and the code
rejects it. Overlapping observation and offset uncertainty produce lower bound
zero; averaging endpoints must not manufacture positive evidence.
\Listing{calibrated_interval}
The assigned example $y\in[.48,.52]$, $o\in[.08,.12]$, $g\in[.9,1.1]$ gives
$B\in[.32727,.48889]$. It does not guarantee $B\ge.45$.
This is a worst-case enclosure, not a confidence or credible interval.
Correlation, saturation, photobleaching and intermediate activation alter the
observation model.
\Plate{06-evidence}{Calibration intervals and independent-unit hierarchy. The unfavorable second unit in this synthetic dataset is retained.}{fig:dna31-evidence}

\section{Count preparations, not repeated readings}
Three independently prepared samples read twice are three independent units,
not six. Technical readings $(.40,.42)$, $(.10,.12)$ and $(.38,.40)$ give
unit means $.41$, $.11$, $.39$ and equal-unit mean $.30333$.
The low second unit is not dropped for disagreeing with an intended narrative.
\Listing{summarize_units}
With one independent unit, between-unit sample standard deviation is unidentified:
the function returns \texttt{None}, not zero. Descriptive sample spread is still
not an interval for future samples. Batch effects, dependence, sampling design
and exclusion rules require separate analysis.

\section{Evidence scopes are not a ranking}
A formal theorem, simulated trajectory, buffer measurement and cellular
measurement answer different questions. An organism-level observation is not a
stronger theorem, nor does it automatically validate a buffer-specific kinetic
claim. The incoming ordinal environment gate is replaced by exact kind, context
and protocol matching.
\Listing{admissible}
The \Term{evidence coverage gate} checks controls, observables, distinct unit
identifiers and a source locator. It does not read raw values, establish an
effect, or perform scientific peer review. A negative experiment may satisfy
metadata coverage perfectly.

Transporting a claim requires an explicit bridge: what mechanism is shared,
which rates change, which delivery barrier appears, and what remains
identifiable? Such a bridge is scientific work, not an ordering comparison
between environment labels.

\section{Make claims falsifiable}
One supported claim is: under this mass-action model and the sufficient Euler
bound, updates preserve nonnegative concentrations and the two pools up to
roundoff. A valid state satisfying the bound but violating positivity or
systematically drifting a pool would falsify that implementation claim.

A biological claim would instead name a reporter, environment, time window,
output criterion and control panel. It requires raw data, calibrated observations
and independent units. Numerical checks cannot establish that claim.
This distinction keeps an implementation proof from becoming an unearned
statement of cellular capability.

\section{Worked exercises}
\begin{enumerate}
\item \textbf{Units.} With micromolar and seconds, what units does $k_{\rm on}$
require? \emph{Answer:} $\mu\mathrm M^{-1}\mathrm s^{-1}$.
\item \textbf{Pool proof.} Apply $(1,0,1,1,0)$ to each column of $S$.
\emph{Answer:} all results are zero, hence $\dot C_A=0$.
\item \textbf{Capacity.} Can $B=.2$ occur from $A_0=1,F_0=.1$ and empty
other pools? \emph{Answer:} no; $B=W\le C_F=.1$.
\item \textbf{Projection.} Is $A$ all unconverted input?
\emph{Answer:} no; captured $I$ is also unconverted, giving $A+I$.
\item \textbf{Initial layer.} Compare exact and reduced $\dot B(0)$.
\emph{Answer:} zero versus $k_{\rm eff}A_0$.
\item \textbf{Repeated pole.} For $u=k_{\rm cat}=1,k_{\rm off}=0$, find
$I(1)/A_0$. \emph{Answer:} $e^{-1}$, using the limit solution.
\item \textbf{No conversion.} Set $k_{\rm cat}=0$.
\emph{Answer:} $B$ remains unchanged although $I$ may accumulate.
\item \textbf{Unsafe step.} With $A=F=1$, $k_{\rm on}=1$,
$k_{\rm off}=.9,k_{\rm cat}=.1$, assess $h=2$.
\emph{Answer:} $h\max(1,1,1)=2>1$; reject it.
\item \textbf{Clipping.} Why not clamp negative $A$?
\emph{Answer:} it changes the pool without a reaction and hides step failure.
\item \textbf{Order.} Halving a step halves its error against an independent
reference. \emph{Answer:} this indicates first-order behavior in that regime,
not a universal accuracy guarantee.
\item \textbf{Loading.} Find branch-1 probability for hazards $2$ and $3$.
\emph{Answer:} $2/5$, not the isolated value one.
\item \textbf{Zero-$L$ bound.} Use $E_0=.1,\eta=.03,t=2$.
\emph{Answer:} $.16$.
\item \textbf{Certification.} Does evaluating the envelope certify $\eta$?
\emph{Answer:} no; its bound must be justified on the intended domain.
\item \textbf{Ambiguity.} Give two states with $B+I=.5$.
\emph{Answer:} $(.4,.1)$ and $(.2,.3)$.
\item \textbf{Extra channel.} Given $I=.1$ and $B+I=.5$, find $B$.
\emph{Answer:} $.4$ under the linear observation model.
\item \textbf{Threshold.} Does the worked interval guarantee $B\ge.45$?
\emph{Answer:} no; its lower endpoint is about $.3273$.
\item \textbf{Inconsistency.} Can $y\in[0,.1]$, $o\in[.2,.3]$ with positive
gain imply nonnegative $B$? \emph{Answer:} no; reject the assumed model.
\item \textbf{Replication.} One preparation is read 100 times.
\emph{Answer:} one independent unit; between-unit spread is unidentified.
\item \textbf{Transport.} Does a cellular result verify a buffer-A claim?
\emph{Answer:} not automatically; the contexts and relevant mechanisms differ.
\item \textbf{Negative study.} Can a failed experiment pass coverage?
\emph{Answer:} yes; coverage is not a positive-outcome label.
\end{enumerate}

\section{Synthesis and the next chapter}
Formal specification, physical realization, observation and scoped claim are
four coupled contracts. Conservation and exact transients make realization
testable; calibration and unit accounting make its evidence interpretable.
None substitutes for the others.

Chapter 32's molecular/genomic synthesis must preserve these distinctions.
A metaphor between genomic regulation and software is not a kinetic refinement
proof; a simulator output is not a cellular measurement. The next integration
must separate established results, conditional conclusions and research proposals.
